\chapter{Energ\'ia libre hologr\'afico-modular y memoria informacional de la geometr\'ia}
\label{chap:energia-libre-holografica}

\section{Estado bicapa orientado \(T_{A,y}\), normalización \(\rho_y\) y separación entre intensidad y orientación de hoja}

La respuesta angular construida en los cap\'itulos precedentes contiene dos
datos espectrales de naturaleza complementaria. La coordenada
\(A_{\rm rad}\) fija la intensidad com\'un de las dos hojas; la coordenada
orientada
\[
 y=C^{*}_{\rm rad}
\]
fija su separaci\'on. Sean \(P_+\) y \(P_-\) proyecciones ortogonales de
rango uno sobre un espacio de Hilbert bidimensional \(\mathcal H_{\rm bi}\),
y definamos
\begin{equation}
 I_{\rm bi}=P_++P_-,
 \qquad
 R=P_+-P_-,
 \qquad
 R^2=I_{\rm bi}.
 \label{eq:rm-info-graduacion}
\end{equation}
El operador angular positivo es
\begin{equation}
 T_{A,y}
 =\exp(-A_{\rm rad}I_{\rm bi}-yR)
 =q_+P_++q_-P_-,
 \qquad
 q_\pm=\exp(-A_{\rm rad}\mp y).
 \label{eq:rm-info-operador-angular}
\end{equation}
Su normalizaci\'on elimina exactamente la intensidad escalar y conserva la
orientaci\'on relativa:
\begin{equation}
 \boxed{
 \rho_y
 =\frac{T_{A,y}}{\Tr T_{A,y}}
 =\frac{\exp(-yR)}{2\cosh y}
 =p_+(y)P_++p_-(y)P_- ,}
 \label{eq:rm-info-rho-y}
\end{equation}
donde
\begin{equation}
 p_+(y)=\frac{\ee^{-y}}{2\cosh y},
 \qquad
 p_-(y)=\frac{\ee^{y}}{2\cosh y}.
 \label{eq:rm-info-probabilidades-hoja}
\end{equation}
La intensidad \(A_{\rm rad}\) permanece en \(\Tr T_{A,y}\) y en el
determinante del operador no normalizado; el estado \(\rho_y\) representa el
sector informacional puro de la elecci\'on de hoja. Esta separaci\'on tipa dos
lectores del mismo antecedente y evita atribuir a una probabilidad
normalizada la escala constitutiva que fue retirada por la normalizaci\'on.

\begin{definition}[Hamiltoniano hologr\'afico-modular bicapa]
El Hamiltoniano hologr\'afico-modular finito del estado orientado es
\begin{equation}
 \boxed{
 K_y=-\log\rho_y
 =\log(2\cosh y)I_{\rm bi}+yR.}
 \label{eq:rm-info-hamiltoniano-bicapa}
\end{equation}
El t\'ermino escalar normaliza el estado y el t\'ermino \(yR\) conserva la
procedencia de hoja.
\end{definition}

La denominaci\'on \emph{hologr\'afico-modular} declara una funci\'on precisa:
el observable total publica la ocupaci\'on, mientras el generador modular
retiene la distribuci\'on entre sectores que hizo posible esa ocupaci\'on. En
dimensi\'on finita, \(K_y\) es una matriz autoadjunta bien definida. Su
extensi\'on al sistema infinito se formular\'a mediante una red local y la
representaci\'on GNS; ambas residencias se mantienen distinguidas.

\section{Termodinámica espectral exacta de las \(2\) hojas}

La funci\'on de partici\'on reducida y el momento orientado son
\begin{equation}
 Z_{\rm bi}(y)=2\cosh y,
 \qquad
 \langle R\rangle_y=\Tr(\rho_yR)=-\tanh y.
 \label{eq:rm-info-particion-momento}
\end{equation}
La entrop\'ia de von Neumann se calcula sin aproximaci\'on:
\begin{equation}
 \boxed{
 S(\rho_y)
 =-\Tr(\rho_y\log\rho_y)
 =\log(2\cosh y)-y\tanh y.}
 \label{eq:rm-info-entropia-bicapa}
\end{equation}
Su paridad y su variaci\'on satisfacen
\begin{equation}
 S(\rho_{-y})=S(\rho_y),
 \qquad
 \frac{\dd}{\dd y}S(\rho_y)=-y\,\operatorname{sech}^2y.
 \label{eq:rm-info-paridad-entropia}
\end{equation}
El d\'eficit respecto del estado bicapa uniforme queda determinado por
\begin{equation}
 I_{\rm or}^{(2)}(y)
 =\log2-S(\rho_y)
 =y\tanh y-\log\cosh y
 \geq0.
 \label{eq:rm-info-deficit-orientado}
\end{equation}
La igualdad ocurre s\'olo en \(y=0\). El d\'eficit mide la informaci\'on
retenida por la polarizaci\'on de las hojas, mientras la entrop\'ia mide la
distribuci\'on probabil\'istica publicada.

\begin{theorem}[Entrop\'ia relativa de dos orientaciones]
\label{thm:rm-info-relative-general}
Para \(y,z\in\mathbb R\), los estados fieles de
\cref{eq:rm-info-rho-y} satisfacen
\begin{equation}
 \boxed{
 D(\rho_y\Vert\rho_z)
 =\log\frac{\cosh z}{\cosh y}
 +(y-z)\tanh y.}
 \label{eq:rm-info-relative-general}
\end{equation}
En particular, la inversi\'on de hoja posee el coste exacto
\begin{equation}
 \boxed{
 D(\rho_y\Vert\rho_{-y})=2y\tanh y.}
 \label{eq:rm-info-relative-inversion}
\end{equation}
\end{theorem}

\begin{proof}
De \(\log\rho_y=-\log(2\cosh y)I_{\rm bi}-yR\) se sigue
\[
 \log\rho_y-\log\rho_z
 =\log\frac{\cosh z}{\cosh y}I_{\rm bi}+(z-y)R.
\]
Al tomar la expectativa en \(\rho_y\) y usar
\(\langle R\rangle_y=-\tanh y\), se obtiene
\cref{eq:rm-info-relative-general}. Sustituir \(z=-y\) elimina el cociente
de cosenos hiperb\'olicos y produce
\cref{eq:rm-info-relative-inversion}. La expresi\'on es positiva porque
\(y\) y \(\tanh y\) tienen el mismo signo.
\end{proof}

El coste simetrizado de Jeffreys es
\begin{equation}
 D(\rho_y\Vert\rho_{-y})
 +D(\rho_{-y}\Vert\rho_y)
 =4y\tanh y.
 \label{eq:rm-info-jeffreys}
\end{equation}
La curvatura local de la familia exponencial est\'a dada por
\begin{equation}
 \frac{\dd^2}{\dd y^2}\log Z_{\rm bi}(y)
 =\operatorname{Var}_{\rho_y}(R)
 =\operatorname{sech}^2y.
 \label{eq:rm-info-fisher}
\end{equation}
As\'i, la misma coordenada orientada determina tres cantidades con paridades
diferentes: \(\langle R\rangle_y\) es impar; \(S(\rho_y)\) es par; y la
distancia entre orientaciones es positiva y par como funci\'on del m\'odulo
de \(y\). La geometr\'ia, la informaci\'on y la inversi\'on conservan
componentes distintas de un mismo estado.

\begin{corollary}[Separaci\'on de intensidad y orientaci\'on]
El determinante
\(\det T_{A,y}=\exp(-2A_{\rm rad})\) depende de la intensidad angular,
mientras el cociente \(q_-/q_+=\exp(2y)\), el momento
\(\langle R\rangle_y\) y
\(D(\rho_y\Vert\rho_{-y})\) dependen de la orientaci\'on. Por tanto, el
lector par que conduce a \(D_A\) y el lector impar que conduce a la memoria
de hoja son reconstruibles conjuntamente y permanecen tipados.
\end{corollary}

\section{Purificaci\'on termocampo doble y conjugaci\'on de hojas}

Sea \((|+\rangle,|-\rangle)\) una base adaptada a
\((P_+,P_-)\), y sea \(\overline{\mathcal H}_{\rm bi}\) una copia
declarada. Definimos
\begin{equation}
 |\operatorname{TFD}_y\rangle
 =\sqrt{p_+(y)}\,|+\rangle|\bar +\rangle
 +\sqrt{p_-(y)}\,|-\rangle|\bar -\rangle.
 \label{eq:rm-info-tfd-bicapa}
\end{equation}

\begin{proposition}[Purificaci\'on bicapa]
\label{prop:rm-info-tfd-bicapa}
El vector de \cref{eq:rm-info-tfd-bicapa} est\'a normalizado y satisface
\begin{equation}
 \Tr_{\overline{\mathcal H}_{\rm bi}}
 |\operatorname{TFD}_y\rangle
 \langle\operatorname{TFD}_y|
 =\rho_y.
 \label{eq:rm-info-tfd-partial}
\end{equation}
La permutaci\'on simult\'anea de las etiquetas \(+\leftrightarrow-\)
transporta \(|\operatorname{TFD}_y\rangle\) en
\(|\operatorname{TFD}_{-y}\rangle\).
\end{proposition}

\begin{proof}
La suma de los cuadrados de los coeficientes es
\(p_+(y)+p_-(y)=1\). La traza parcial elimina los t\'erminos cruzados y
deja \(p_+(y)P_++p_-(y)P_-\). El intercambio de las etiquetas permuta
ambos coeficientes, pues \(p_+(-y)=p_-(y)\).
\end{proof}

La purificaci\'on demuestra una igualdad de entrop\'ias dentro del sistema
duplicado declarado:
\begin{equation}
 S(\rho_y)
 =S\!\left(
 \Tr_{\mathcal H_{\rm bi}}
 |\operatorname{TFD}_y\rangle
 \langle\operatorname{TFD}_y|
 \right).
 \label{eq:rm-info-tfd-entropy}
\end{equation}
La copia termocampo doble posee aqu\'i una funci\'on matem\'atica: exhibe un
estado puro cuya reducci\'on produce exactamente el estado bicapa. Su
identificaci\'on con dos regiones f\'isicas de un espacio-tiempo requiere
una factorizaci\'on de observables y una interfaz geom\'etrica adicional.

La TFD bicapa y la TFD qutrit del
\cref{thm:modular-tfd} son construcciones relacionadas y distintas. La
primera purifica la orientaci\'on de dos hojas; la segunda purifica los
treinta y seis enlaces qutrit de los canales \(90/120\). La confluencia
entre ambas reside en el antecedente angular y en la etiqueta de
procedencia, no en una identificaci\'on de sus espacios de Hilbert.

\section{Geometr\'ia total y procedencia modular}

Los operadores n\'umero \(N_{90}\) y \(N_{120}\) conmutan. La geometr\'ia
areal publica la suma de ocupaciones,
\begin{equation}
 \widehat{\mathcal A}
 =\gammaH a_0\ell_P^2(N_{90}+N_{120}),
 \label{eq:rm-info-area-total}
\end{equation}
mientras el Hamiltoniano modular de borde conserva su procedencia,
\begin{equation}
 K_A-cI
 =\Dmod(90N_{90}+120N_{120}).
 \label{eq:rm-info-hhm-provenance}
\end{equation}
Definamos los observables adimensionales
\begin{equation}
 N=N_{90}+N_{120},
 \qquad
 D=90N_{90}+120N_{120}.
 \label{eq:rm-info-ND}
\end{equation}

\begin{theorem}[Reconstrucci\'on conjunta de los canales]
\label{thm:rm-info-reconstruction-90-120}
El par \((N,D)\), y por tanto el par
\((\widehat{\mathcal A},K_A)\) una vez fijadas sus escalas, reconstruye
exactamente los dos operadores de ocupaci\'on:
\begin{equation}
 \boxed{
 N_{90}=4N-\frac{D}{30},
 \qquad
 N_{120}=\frac{D}{30}-3N.}
 \label{eq:rm-info-inverse-90-120}
\end{equation}
\end{theorem}

\begin{proof}
El sistema lineal posee matriz
\(\left(\begin{smallmatrix}1&1\\90&120\end{smallmatrix}\right)\)
de determinante \(30\). Su inversa da
\cref{eq:rm-info-inverse-90-120}. La conmutatividad de los operadores
n\'umero permite aplicar la identidad en su c\'alculo funcional conjunto.
\end{proof}

El resultado formula una holograf\'ia de reconstrucci\'on concreta. El
\'area determina la cantidad geom\'etrica total; el Hamiltoniano modular
determina la procedencia ponderada; la pareja recupera la distribuci\'on
sectorial. Dos estados con la misma \`area pueden portar historias
distintas, y esa diferencia permanece visible en \(K_A\). El coeficiente
\(30\) que hace invertible el sistema es el mismo defecto de completaci\'on
entre los canales \(90\) y \(120\).

La reconstrucci\'on exige conservar ambos observables. Una compresi\'on que
retuviera s\'olo \(N\) identificar\'ia todos los pares
\((N_{90},N_{120})\) de suma constante. Una compresi\'on que retuviera s\'olo
\(D\) identificar\'ia pares de igual peso modular. El par conjunto posee
rango dos y conserva la genealog\'ia sectorial pertinente.

\section{1.ª ley local y balance finito de entropía relativa}

La formulaci\'on finita emplea la entrop\'ia relativa de estados operatorios
en el sentido de Araki \cite{Araki1976}.

Sea \(\rho_A\) el estado qutrit de borde de
\cref{eq:modular-product-state} y sea
\(K_A=-\log\rho_A\). La descomposici\'on areal demostrada en el
\cref{thm:modular-first-law} adopta la forma
\begin{equation}
 K_A-cI
 =\beta_{90}\mathcal A_{90}
 +\beta_{120}\mathcal A_{120},
 \qquad
 \frac{\beta_{120}}{\beta_{90}}=\frac43.
 \label{eq:rm-info-K-area}
\end{equation}
Para una curva diferenciable de estados normalizados que atraviesa
\(\rho_A\), la variaci\'on tangente satisface
\begin{equation}
 \boxed{
 \delta S
 =\beta_{90}\,\delta\langle\mathcal A_{90}\rangle
 +\beta_{120}\,\delta\langle\mathcal A_{120}\rangle.}
 \label{eq:rm-info-first-law}
\end{equation}
El cociente \(4/3\) registra el distinto peso informacional de los dos
sectores; un mismo incremento de \`area total puede producir variaciones
entr\'opicas diferentes seg\'un su canal de procedencia.

La formulaci\'on finita conserva el t\'ermino que la linealizaci\'on elimina.
Si \(\sigma\) es un estado cuyo soporte est\'a contenido en el de
\(\rho_A\), entonces
\begin{equation}
 \boxed{
 \Delta S
 =\beta_{90}\Delta\langle\mathcal A_{90}\rangle
 +\beta_{120}\Delta\langle\mathcal A_{120}\rangle
 -D(\sigma\Vert\rho_A).}
 \label{eq:rm-info-finite-law}
\end{equation}

\begin{proof}
La definici\'on de entrop\'ia relativa proporciona
\[
 D(\sigma\Vert\rho_A)
 =-S(\sigma)+\Tr(\sigma K_A).
\]
Para \(\rho_A\), la misma expresi\'on vale cero. Al restar ambas
identidades se obtiene
\(\Delta S=\Delta\langle K_A\rangle-D(\sigma\Vert\rho_A)\).
El t\'ermino escalar de \(K_A\) se cancela entre estados normalizados, y
\cref{eq:rm-info-K-area} concluye la prueba.
\end{proof}

La ley local expresa la derivada de la entrop\'ia en el estado de
referencia. La ley finita expresa el balance exacto entre respuesta
geom\'etrica, cambio entr\'opico y distancia informacional. La positividad de
\(D\) cuantifica la convexidad que separa el desplazamiento finito de su
aproximaci\'on tangente.

\section{Energ\'ia libre hologr\'afico-modular de horizonte}

Consideremos una realizaci\'on de horizonte en la que un radio areal \(R\)
determina
\begin{equation}
 E_\Lambda(R)=\frac{c^4R}{2G},
 \qquad
 k_BT_H(R)=\frac{\hbar c}{2\pi R},
 \qquad
 \frac{S_H(R)}{k_B}=\frac{\pi R^2}{\ell_P^2}.
 \label{eq:rm-info-horizon-data}
\end{equation}
La identidad \(\ell_P^2=G\hbar/c^3\) da
\begin{equation}
 \boxed{E_\Lambda(R)=T_H(R)S_H(R).}
 \label{eq:rm-info-horizon-saturation}
\end{equation}
Para una perturbaci\'on finita \(\sigma\) del estado de referencia
\(\rho_A\), escribamos
\begin{equation}
 \Delta k_A=\Tr((\sigma-\rho_A)K_A),
 \qquad
 \Delta s_A=S(\sigma)-S(\rho_A),
 \label{eq:rm-info-delta-k-s}
\end{equation}
y definamos las magnitudes de respuesta
\begin{equation}
 \mathcal E_R(\sigma)
 =E_\Lambda(R)+k_BT_H(R)\Delta k_A,
 \qquad
 \mathcal S_R(\sigma)
 =S_H(R)+k_B\Delta s_A.
 \label{eq:rm-info-response-energy-entropy}
\end{equation}

\begin{theorem}[Identidad de energ\'ia libre hologr\'afico-modular]
\label{thm:rm-info-free-energy}
Bajo la realizaci\'on de horizonte
\cref{eq:rm-info-horizon-data}, para todo \(\sigma\) con soporte contenido
en el de \(\rho_A\),
\begin{equation}
 \boxed{
 \mathcal E_R(\sigma)-T_H(R)\mathcal S_R(\sigma)
 =k_BT_H(R)D(\sigma\Vert\rho_A)
 \geq0.}
 \label{eq:rm-info-free-energy}
\end{equation}
\end{theorem}

\begin{proof}
La saturaci\'on de referencia
\cref{eq:rm-info-horizon-saturation} cancela el t\'ermino
\(E_\Lambda-T_HS_H\). Por definici\'on,
\[
 \mathcal E_R-T_H\mathcal S_R
 =k_BT_H(\Delta k_A-\Delta s_A).
\]
La identidad finita de entrop\'ia relativa da
\(\Delta k_A-\Delta s_A=D(\sigma\Vert\rho_A)\), y la desigualdad de Klein
da la positividad.
\end{proof}

En el cierre \(R=\ell_P\), el cuanto de Planck
\(E_P=\hbar c/\ell_P\) produce
\begin{equation}
 k_BT_H(\ell_P)=\frac{E_P}{2\pi},
 \qquad
 E_\Lambda(\ell_P)=\frac{E_P}{2},
 \label{eq:rm-info-planck-horizon-data}
\end{equation}
y por ello
\begin{equation}
 \boxed{
 \mathcal E_{\ell_P}(\sigma)
 -T_H(\ell_P)\mathcal S_{\ell_P}(\sigma)
 =\frac{E_P}{2\pi}D(\sigma\Vert\rho_A).}
 \label{eq:rm-info-planck-free-energy}
\end{equation}
La identidad estructural
\begin{equation}
 E_P=k_B\Theta_{\rm clk}\log3
 \label{eq:rm-info-planck-trit}
\end{equation}
convierte el coeficiente del defecto en
\begin{equation}
 \frac{E_P}{2\pi}
 =\frac{k_B\Theta_{\rm clk}\log3}{2\pi}.
 \label{eq:rm-info-trit-free-energy}
\end{equation}
El horizonte de referencia satura el balance; la desviaci\'on finita queda
medida por el coste informacional de abandonar su genealog\'ia modular.

La realizaci\'on tambi\'en satisface
\begin{equation}
 E_\Lambda(\ell_P)t_P=\frac{\hbar}{2},
 \qquad
 E_\Lambda(\ell_P)T_{108}=\frac h2,
 \label{eq:rm-info-half-action}
\end{equation}
cuando \(t_P=\ell_P/c\) y \(T_{108}=2\pi t_P\). La media acci\'on
resulta del acoplamiento entre periodicidad t\'ermica y geometr\'ia areal.
Esta igualdad se mantiene separada de la identidad de energ\'ia libre: la
primera fija el cuanto del horizonte de referencia; la segunda cuantifica
las perturbaciones del estado modular.

\section{Confluencia entre cuanto areal y cuanto informacional}

La transici\'on hexada--octada publica el cuanto geom\'etrico
\begin{equation}
 \delta\mathcal A
 =\gammaH a_0\ell_P^2,
 \label{eq:rm-info-area-quantum}
\end{equation}
y el reloj CT108 publica el cuanto informacional
\begin{equation}
 E_P=k_B\Theta_{\rm clk}\log3.
 \label{eq:rm-info-energy-quantum}
\end{equation}
Su cociente define la tensi\'on tipada
\begin{equation}
 \boxed{
 \mathcal T_{I/A}
 =\frac{E_P}{\delta\mathcal A}
 =\frac{k_B\Theta_{\rm clk}\log3}
 {\gammaH a_0\ell_P^2}.}
 \label{eq:rm-info-tension-information-area}
\end{equation}
La ecuaci\'on asigna una conversi\'on entre una decisi\'on ternaria y una
unidad de geometr\'ia. El numerador pertenece a la rama
acci\'on--reloj--informaci\'on; el denominador pertenece a la rama
hexada--octada--\`area. Ambas descienden del mismo estado enriquecido y
confluyen en gravedad despu\'es de conservar sus dominios.

El par
\begin{equation}
 \left(
 \widehat{\mathcal A},K_A
 \right)
 \label{eq:rm-info-area-K-pair}
\end{equation}
reconstruye cantidad y procedencia; el par
\begin{equation}
 \left(
 \delta\mathcal A,E_P
 \right)
 \label{eq:rm-info-area-energy-pair}
\end{equation}
relaciona geometr\'ia elemental y coste informacional. La primera pareja es
espectral; la segunda es dimensional. Juntas describen c\'omo una estructura
discreta de orientaci\'on se publica como \`area, energ\'ia y entrop\'ia.

\section{Red local UHF y generador modular en la representaci\'on GNS}

Sea \((\mathcal A_m,\iota_m)\) una red de \`algebras matriciales locales
con
\begin{equation}
 \mathcal A_m\simeq M_{9^m}(\mathbb C),
 \qquad
 \iota_m(a)=a\otimes I_9,
 \qquad
 \mathcal A_{\rm UHF}=\varinjlim_m\mathcal A_m.
 \label{eq:rm-info-uhf-net}
\end{equation}
Una familia de estados locales \(\omega_m\) define un estado global
\(\omega\) cuando satisface
\begin{equation}
 \omega_{m+1}\circ\iota_m=\omega_m.
 \label{eq:rm-info-state-compatibility}
\end{equation}
En cada corte finito puede escribirse
\begin{equation}
 \omega_m(a)=\Tr(\rho_m a),
 \qquad
 K_m=-\log\rho_m.
 \label{eq:rm-info-local-densities}
\end{equation}
La familia \((K_m)\) constituye la residencia local del Hamiltoniano
hologr\'afico-modular.

En la representaci\'on GNS
\((\pi_\omega,\mathcal H_\omega,\Omega_\omega)\), la teor\'ia modular se
expresa mediante el operador modular \(\Delta_\omega\) y el grupo
\begin{equation}
 \sigma_t^\omega(x)
 =\Delta_\omega^{\,it}x\Delta_\omega^{-it},
 \qquad
 x\in\pi_\omega(\mathcal A_{\rm UHF})''.
 \label{eq:rm-info-gns-modular-flow}
\end{equation}
El generador \(-\log\Delta_\omega\) puede ser no acotado y pertenece a la
representaci\'on est\'andar; una matriz densidad global traza-clase dentro de
la UHF no forma parte, en general, de esta construcci\'on.

La expectativa tracial non\'adica
\begin{equation}
 E_m=\operatorname{id}\otimes\tau_9:
 \mathcal A_{m+1}\longrightarrow\mathcal A_m
 \label{eq:rm-info-tracial-expectation}
\end{equation}
preserva la unidad y constituye el lector can\'onico de compresi\'on
matricial. La compatibilidad de un estado modular exige adem\'as
\begin{equation}
 \omega_m\circ E_m=\omega_{m+1}
 \quad\hbox{en el subespacio declarado,}
 \label{eq:rm-info-expectation-state}
\end{equation}
o la condici\'on equivalente apropiada a la inclusi\'on utilizada. La
tracialidad del lector no convierte autom\'aticamente un producto Gibbs no
tracial en una familia compatible.

El \cref{thm:modular-typeIII} demuestra, bajo persistencia fiel de los dos
pesos \(90\Dmod\) y \(120\Dmod\), densidad asint\'otica positiva de ambos
canales, factorialidad y ausencia de cocientes adicionales, que el grupo de
razones es
\begin{equation}
 90\Dmod\mathbb Z+120\Dmod\mathbb Z
 =30\Dmod\mathbb Z
 \label{eq:rm-info-ratio-group}
\end{equation}
y que el cierre ITPFI hiperinfinito es de tipo
\begin{equation}
 \boxed{
 III_{\lambda_\partial},
 \qquad
 \lambda_\partial=\exp(-30\Dmod).}
 \label{eq:rm-info-type-iii}
\end{equation}
El enunciado clasifica la red producto especificada. Una modificaci\'on del
grupo de razones, la desaparici\'on asint\'otica de un canal o la p\'erdida de
factorialidad obliga a recalcular el tipo. En particular, el cierre tracial
II\(_1\) de la UHF y el cierre modular tipo III responden a estados y
representaciones diferentes; ambos son lectores operatorios del sistema
non\'adico y ninguno sustituye al antecedente APP--TRIT--TPK.

\section{Identidades exactas del cierre gravitatorio y condiciones de la realización tensorial}

La arquitectura de este cap\'itulo distingue los resultados siguientes:
\begin{enumerate}[label=\textup{(\roman*)}]
 \item La normalizaci\'on bicapa, el Hamiltoniano
 \(K_y\), la entrop\'ia
 \(S(\rho_y)\) y la identidad
 \(D(\rho_y\Vert\rho_{-y})=2y\tanh y\) son
 \status{formalizaci\'on nueva exacta} de la respuesta angular ya
 construida.
 \item La TFD bicapa es una \status{formalizaci\'on nueva exacta} en el
 espacio duplicado declarado. La interpretaci\'on de la copia como otra
 regi\'on f\'isica requiere una factorizaci\'on geom\'etrica posterior.
 \item La reconstrucci\'on \(90/120\), la TFD qutrit, la primera ley local,
 la identidad finita y la clasificaci\'on modular bajo sus hip\'otesis son
 \status{resultados exactos} de la construcción operatorial precedente.
 \item La igualdad de horizonte
 \(E_\Lambda=T_HS_H\) es exacta en la carta de Sitter declarada. Su
 composici\'on con el estado qutrit para formar
 \cref{eq:rm-info-free-energy} es una
 \status{interfaz f\'isica condicionada} a esa realizaci\'on.
 \item El entrelazador entre microestados de horizonte, TFD de borde y
 autoespacios hiperb\'olicos queda clasificado por espectro y multiplicidad
 en el \cref{thm:ivv-horizon-tfd-perron-classification}: la coincidencia
 tipada produce la equivalencia unitaria y su fallo demuestra la obstrucci\'on.
\end{enumerate}

\begin{proofstatusbox}
Las identidades espectrales de la familia \(\rho_y\), su purificaci\'on,
la reconstrucci\'on lineal de los sectores \(90/120\) y el balance finito
con entrop\'ia relativa son exactos en dimensi\'on finita. La clasificaci\'on
tipo III es exacta bajo las cuatro hip\'otesis formuladas en esta construcción.
La energ\'ia libre gravitatoria se afirma en la realizaci\'on de horizonte de
\cref{eq:rm-info-horizon-data}; conserva la residencia propia del estado
modular y mantiene la din\'amica cosmol\'ogica como una promoci\'on posterior.
\end{proofstatusbox}

\begin{falsifierbox}
Refutan las afirmaciones correspondientes: un espectro de \(\rho_y\) que no
sume uno; un c\'alculo de
\(D(\rho_y\Vert\rho_{-y})\) distinto de \(2y\tanh y\); una traza parcial de
la TFD que no produzca \(\rho_y\); dos pares distintos
\((N_{90},N_{120})\) con los mismos \((N,D)\); una perturbaci\'on finita
que viole \cref{eq:rm-info-finite-law}; una carta de horizonte que incumpla
\(E_\Lambda=T_HS_H\); o una red infinita cuyo grupo de razones difiera de
\(30\Dmod\mathbb Z\) bajo las hip\'otesis declaradas. Tambi\'en refuta una
promoci\'on operatoria la sustituci\'on de la red local por una matriz densidad
global inexistente o la identificaci\'on del cierre tracial II\(_1\) con el
cierre modular tipo III.
\end{falsifierbox}
